%% ============================================================
%% PARTE I — FUNDAMENTAÇÃO TEÓRICA
%% ============================================================

\chapter{O Método Híbrido Fonico-Kumon na Alfabetização}

\section{O que é o Método Híbrido?}

O método híbrido Fonico-Kumon é uma abordagem pedagógica que integra dois pilares fundamentais para a alfabetização de crianças do 1º ao 5º ano do Ensino Fundamental.

\subsection{O Método Fonético (Phonics)}

O método fonético ensina que as letras do alfabeto representam sons (fonemas). Esta abordagem possui base científica sólida, com pesquisas como a de Ehri (2005) demonstrando que o ensino sistemático de fonemas é o mais eficaz para alfabetização.

\textbf{Princípios do Método Fonético:}
\begin{enumerate}[leftmargin=2cm]
  \item \textbf{Princípio alfabético}: letras representam sons
  \item \textbf{Progressão sistemática}: fonema $\rightarrow$ grafema $\rightarrow$ sílaba $\rightarrow$ palavra $\rightarrow$ frase
  \item \textbf{Decodificação}: capacidade de ``traduzir'' palavras desconhecidas
  \item \textbf{Consciência fonológica}: percepção dos sons da língua
\end{enumerate}

\subsection{O Método Kumon}

Desenvolvido pelo educador japonês Toru Kumon em 1954, o método Kumon é baseado em três princípios fundamentais:

\begin{enumerate}[leftmargin=2cm]
  \item \textbf{Mastery-based} (domínio antes de avançar): o aluno só avança quando domina o conteúdo atual
  \item \textbf{Progressão granular}: cada habilidade é decomposta em micro-passos incrementais
  \item \textbf{Prática diária}: folhas de trabalho estruturadas com exercícios progressivos
\end{enumerate}

\subsection{Por que um Método Híbrido?}

A integração dos dois métodos permite combinar o melhor de cada abordagem:

\begin{table}[h]
\centering
\caption{Comparação entre Métodos}
\begin{tabular}{|l|c|c|c|}
\hline
\textbf{Característica} & \textbf{Fonético} & \textbf{Kumon} & \textbf{Híbrido} \\
\hline
Progressão & Sistemática & Granular & Sistemática + Granular \\
Prática & Em sala & Diária (folhas) & Em sala + em casa \\
Avaliação & Diagnóstica & Mastery-based & Diagnóstica + Mastery \\
Flexibilidade & Média & Baixa & Alta \\
Engajamento & Médio & Alto & Alto \\
\hline
\end{tabular}
\end{table}

\section{Princípios Fundamentais}

\begin{enumerate}[leftmargin=2cm]
  \item \textbf{Decomposição}: Cada habilidade de leitura/escrita é decomposta em micro-habilidades
  \item \textbf{Progressão}: De simples para complexo, sem saltos
  \item \textbf{Repetição com variação}: O mesmo conteúdo é praticado com variações incrementais
  \item \textbf{Domínio antes de avançar}: Não se avança sem dominar o nível atual
  \item \textbf{Multisensorial}: Visão + audição + movimento + fala
\end{enumerate}

%% ============================================================

\chapter{A Técnica da Boquinha}

\section{Conceito}

A ``Boquinha'' é um recurso didático multissensorial que associa quatro dimensões de aprendizagem:

\begin{enumerate}[leftmargin=2cm]
  \item \textbf{Visual}: A forma da letra e a imagem da boca
  \item \textbf{Auditiva}: O som (fonema) que a letra representa
  \item \textbf{Motor-oral}: A posição da boca ao pronunciar o fonema
  \item \textbf{Motor-gráfico}: O traçado da letra com direção do movimento
\end{enumerate}

\section{Fundamentação Teórica}

A técnica baseia-se em princípios de neurociência da aprendizagem:

\begin{enumerate}[leftmargin=2cm]
  \item \textbf{Codificação dual} (Paivio, 1986): Informações processadas tanto visual quanto verbalmente são melhor armazenadas
  \item \textbf{Memória procedural}: A associação boca-letra cria ``rotinas motoras'' de fixação
  \item \textbf{Consciência fonológica}: O aluno aprende a articular conscientemente cada som
  \item \textbf{Transferência intermodal}: A posição da boca ``ancora'' a forma da letra na memória
\end{enumerate}

\section{Categorias de Boquinhas}

\subsection{Grupo 1 --- Vogais (boca aberta)}

\begin{boquinha}[Boquinha da Letra A --- Boca Aberta]
\begin{minipage}{0.45\textwidth}
\centering
\textbf{Posição da Boca}\\[0.5em]
Boca bem ABERTA\\
Ar sai forte pela boca\\
Som: ``ÁÁÁÁÁÁÁÁ''
\end{minipage}
\hfill
\begin{minipage}{0.45\textwidth}
\centering
\textbf{Traçado da Letra}\\[0.5em]
\letra{A}
Traçado:
\begin{enumerate}
  \item Diagonal esquerda ($\swarrow$)
  \item Diagonal direita ($\searrow$)
  \item Horizontal ($\rightarrow$)
\end{enumerate}
\end{minipage}
\end{boquinha}

\vspace{1em}

\begin{boquinha}[Boquinha da Letra E --- Boca Entreaberta]
\begin{minipage}{0.45\textwidth}
\centering
\textbf{Posição da Boca}\\[0.5em]
Boca ENTREABERTA\\
Ar sai suave\\
Som: ``ÉÉÉÉÉÉÉÉ''
\end{minipage}
\hfill
\begin{minipage}{0.45\textwidth}
\centering
\textbf{Traçado da Letra}\\[0.5em]
\letra{E}
Traçado:
\begin{enumerate}
  \item Vertical ($\downarrow$)
  \item Horizontal ($\rightarrow$)
\end{enumerate}
\end{minipage}
\end{boquinha}

\vspace{1em}

\begin{boquinha}[Boquinha da Letra I --- Boca Entreaberta]
\begin{minipage}{0.45\textwidth}
\centering
\textbf{Posição da Boca}\\[0.5em]
Boca ENTREABERTA\\
Canto da boca para TRÁS\\
Som: ``III-II-III''
\end{minipage}
\hfill
\begin{minipage}{0.45\textwidth}
\centering
\textbf{Traçado da Letra}\\[0.5em]
\letra{I}
Traçado:
\begin{enumerate}
  \item Vertical ($\downarrow$)
\end{enumerate}
\end{minipage}
\end{boquinha}

\vspace{1em}

\begin{boquinha}[Boquinha da Letra O --- Boca Arredondada]
\begin{minipage}{0.45\textwidth}
\centering
\textbf{Posição da Boca}\\[0.5em]
Boca ARREDONDADA\\
Ar sai em círculo\\
Som: ``ÓÓÓÓÓÓÓÓ''
\end{minipage}
\hfill
\begin{minipage}{0.45\textwidth}
\centering
\textbf{Traçado da Letra}\\[0.5em]
\letra{O}
Traçado:
\begin{enumerate}
  \item Círculo (hora 12 $\rightarrow$ 3 $\rightarrow$ 6 $\rightarrow$ 9 $\rightarrow$ 12)
\end{enumerate}
\end{minipage}
\end{boquinha}

\vspace{1em}

\begin{boquinha}[Boquinha da Letra U --- Boca Bem Arredondada]
\begin{minipage}{0.45\textwidth}
\centering
\textbf{Posição da Boca}\\[0.5em]
Boca BEM ARREDONDADA\\
Ar sai focado\\
Som: ``UUU-UUU-UUU''
\end{minipage}
\hfill
\begin{minipage}{0.45\textwidth}
\centering
\textbf{Traçado da Letra}\\[0.5em]
\letra{U}
Traçado:
\begin{enumerate}
  \item Vertical ($\downarrow$)
  \item Curva ($\rightarrow$)
  \item Vertical ($\uparrow$)
\end{enumerate}
\end{minipage}
\end{boquinha}

\subsection{Grupo 2 --- Consonantes Bilabiais (boca fecha)}

\begin{boquinha}[Boquinha da Letra M --- Boca Fechada, Nariz Vibra]
\begin{minipage}{0.45\textwidth}
\centering
\textbf{Posição da Boca}\\[0.5em]
Boca FECHADA\\
Nariz VIBRA (son nasal)\\
Som: ``MMM-MMM-MMM''
\end{minipage}
\hfill
\begin{minipage}{0.45\textwidth}
\centering
\textbf{Traçado da Letra}\\[0.5em]
\letra{M}
Traçado:
\begin{enumerate}
  \item Vertical ($\downarrow$)
  \item Diagonal ($\searrow$)
  \item Diagonal ($\nearrow$)
  \item Vertical ($\downarrow$)
\end{enumerate}
\end{minipage}
\end{boquinha}

\subsection{Grupo 3 --- Consonantes Linguais (língua toca o céu da boca)}

\begin{boquinha}[Boquinha da Letra T --- Língua nos Dentes]
\begin{minipage}{0.45\textwidth}
\centering
\textbf{Posição da Boca}\\[0.5em]
Boca ENTREABERTA\\
Língua toca atrás dos dentes superiores\\
Som: ``T-T-T'' (seco, rápido)
\end{minipage}
\hfill
\begin{minipage}{0.45\textwidth}
\centering
\textbf{Traçado da Letra}\\[0.5em]
\letra{T}
Traçado:
\begin{enumerate}
  \item Horizontal ($\rightarrow$)
  \item Vertical ($\downarrow$) cruzando
\end{enumerate}
\end{minipage}
\end{boquinha}

%% ============================================================

\chapter{Alinhamento Curricular}

\section{Habilidades de Língua Portuguesa --- 1º Ano (BNCC)}

\begin{table}[h]
\centering
\caption{Habilidades BNCC do 1º Ano e Abordagem no Livro}
\footnotesize
\begin{tabular}{|p{2cm}|p{6cm}|p{4cm}|}
\hline
\textbf{Código} & \textbf{Habilidade} & \textbf{Abordagem} \\
\hline
EF01LP01 & Reconhecer que as letras do alfabeto representam fonemas & Volume 1, Unidades 1--3 \\
\hline
EF01LP02 & Identificar as letras do alfabeto e seus usos & Volume 1, Unidades 1--3 \\
\hline
EF01LP03 & Associar letras e fonemas em palavras & Volume 1, Unidades 4--6 \\
\hline
EF01LP04 & Produzir textos orais e escritos & Volume 1--2, integração \\
\hline
EF01LP05 & Reconhecer a função social da escrita & Volume 1, Parte I \\
\hline
EF01LP06 & Compreender textos lidos pelo professor & Volume 1, Parte II \\
\hline
EF01LP07 & Produzir textos orais e interpretar gêneros & Volume 1, Parte II \\
\hline
EF01LP08 & Ler e escrever palavras e frases simples & Volume 1, Unidades 5--6 \\
\hline
EF01LP09 & Reconhecer e usar palavras comuns & Volume 1, integração \\
\hline
\end{tabular}
\end{table}

\section{Progressão BNCC --- 1º ao 5º Ano}

\begin{center}
\begin{tikzpicture}[
  level/.style={rectangle, draw=primary, fill=primary!10, 
                minimum width=3cm, minimum height=1cm, align=center, font=\small},
  arr/.style={->, thick, color=accent}
]
  %% 1º Ano
  \node[level] (y1) at (0,4) {\textbf{1º Ano}\\Consciência fonológica\\Letras isoladas\\Sílabas diretas};
  
  %% 2º Ano
  \node[level] (y2) at (4,4) {\textbf{2º Ano}\\Sílabas simples\\Palavras simples\\Frases simples};
  
  %% 3º Ano
  \node[level] (y3) at (8,4) {\textbf{3º Ano}\\Sílabas compostas\\Palavras complexas\\Textos curtos};
  
  %% 4º Ano
  \node[level] (y4) at (2,0) {\textbf{4º Ano}\\Textos completos\\Produção escrita\\Análise gramatical};
  
  %% 5º Ano
  \node[level] (y5) at (6,0) {\textbf{5º Ano}\\Textos complexos\\Produção argumentativa\\Revisão e reescrita};
  
  %% Setas
  \draw[arr] (y1) -- (y2);
  \draw[arr] (y2) -- (y3);
  \draw[arr] (y3) -- (y4);
  \draw[arr] (y4) -- (y5);
\end{tikzpicture}
\end{center}

\section{Diretrizes Curriculares Nacionais (DCN)}

As Diretrizes Curriculares Nacionais para o Ensino Fundamental de Nove Anos são estabelecidas pela \textbf{Resolução CNE/CEB nº 7, de 14 de dezembro de 2010}, que fixa as diretrizes curriculares nacionais para o Ensino Fundamental de 9 (nove) anos e complementa a BNCC ao definir os princípios e orientações gerais que devem nortear a organização curricular da educação básica brasileira.

\subsection{Princípios Fundamentais das DCN}

As DCN do Ensino Fundamental (Resolução CNE/CEB nº 7/2010) estabelecem um conjunto de princípios que orientam a prática pedagógica:

\begin{enumerate}[leftmargin=2cm]
  \item \textbf{Éticos}: de justiça, solidariedade, liberdade e autonomia; de respeito à dignidade da pessoa humana e de compromisso com a promoção do bem de todos
  \item \textbf{Políticos}: de reconhecimento dos direitos e deveres de cidadania, de respeito ao bem comum e à preservação do regime democrático e dos recursos ambientais
  \item \textbf{Estéticos}: do cultivo da sensibilidade juntamente com o da racionalidade; do enriquecimento das formas de expressão e do exercício da criatividade
  \item \textbf{Dignidade humana}: respeito à integridade física, afetiva, intelectual e psicológica do educando
  \item \textbf{Formação integral}: desenvolvimento harmonioso de todas as dimensões da pessoa (física, cognitiva, afetiva, ética e estética)
  \item \textbf{Qualidade social da educação}: garantia de aprendizagens significativas para todos os alunos, com respeito à diversidade
  \item \textbf{Inclusão educacional}: atendimento à diversidade de ritmos, condições e necessidades de aprendizagem (Art. 29 e Art. 37 da Res. 7/2010)
\end{enumerate}

\subsection{Linguagens Orais e Escritas como Áreas Obrigatórias}

O Art. 32 da Lei de Diretrizes e Bases da Educação Nacional (Lei nº 9.394/1996), com a redação dada pela Lei nº 11.274/2006, estabelece que o ensino fundamental obrigatório tem por objetivo a formação básica do cidadão mediante o desenvolvimento da capacidade de aprender, tendo como meios básicos o pleno domínio da leitura, da escrita e do cálculo. As DCN reforçam essa obrigatoriedade ao determinar que as línguas orais e escritas constituem uma das áreas de conhecimento que devem compor a estrutura curricular do Ensino Fundamental.

No contexto do método híbrido Fonico-Kumon, essa diretriz é plenamente atendida, pois o livro didático contempla tanto o trabalho com a língua oral (dita, escutada, narrada) quanto a língua escrita (lida, produzida, reescrita), integrando leitura, escrita, oralidade e produção textual de forma articulada.

\subsection{Articulação entre Disciplinas}

O Art. 26 da LDB determina que ``os currículos do ensino fundamental e médio devem ter uma base nacional comum a ser complementada, em cada sistema de ensino e estabelecimento escolar, por uma parte diversificada, exigida pelas características regionais e locais da sociedade, da cultura, da economia e dos educandos''. As DCN aprofundam esse princípio ao orientar que a articulação entre as diferentes disciplinas deve ocorrer de maneira transversal, promovendo a interdisciplinaridade e a contextualização dos conteúdos.

O método híbrido Fonico-Kumon alinha-se a essa diretriz ao integrar:

\begin{itemize}[leftmargin=2cm]
  \item \textbf{Língua Portuguesa}: fonemas, grafemas, ortografia, produção textual
  \item \textbf{Artes}: apreciação estética do traçado, criatividade na escrita
  \item \textbf{Educação Física}: coordenação motora fina e grossa associada ao ato de escrever
  \item \textbf{História e Geografia}: leitura e produção de textos com temas da realidade social
  \item \textbf{Ciências}: compreensão de textos informativos e experimentais
\end{itemize}

\subsection{Resumo das DCN para o Livro Didático}

\begin{table}[h]
\centering
\caption{DCN (Res. CNE/CEB nº 7/2010) e Alinhamento com o Método Híbrido}
\resizebox{\textwidth}{!}{%
\begin{tabular}{|p{3cm}|p{4cm}|p{4cm}|p{4cm}|}
\hline
\textbf{Princípio DCN} & \textbf{O que estabelece} & \textbf{Como o livro atende} & \textbf{Referência} \\
\hline
Dignidade humana & Respeito integral ao educando & Atividades lúdicas e acolhedoras & Art. 12, I \\
\hline
Formação integral & Desenvolvimento multidimensional & Integração oralidade, escrita, arte e movimento & Art. 12, IV \\
\hline
Diversidade & Valorização de culturas diversas & Textos multiculturais e inclusivos & Art. 12, III \\
\hline
Qualidade social & Aprendizagens significativas para todos & Progressão sistemática e mastery-based & Art. 13 \\
\hline
Inclusão educacional & Atendimento à diversidade de ritmos & Rotas 1A--1D e níveis K0--K5 & Art. 29, 37 \\
\hline
Interdisciplinaridade & Articulação entre áreas & Integração Língua Portuguesa com outras disciplinas & Art. 26 (LDB) \\
\hline
\end{tabular}%
}
\end{table}

%% ============================================================
%% NOVO CAPÍTULO: NÍVEIS K INTERNOS (K0-K5)
%% ============================================================

\chapter{Níveis Internos de Progressão (K0--K5)}

\section{Por que uma escala interna própria?}

O método Kumon oficial utiliza uma nomenclatura própria (7A a 2A) que não corresponde diretamente aos anos escolares brasileiros. Para este livro, criamos uma \textbf{escala interna de seis níveis (K0 a K5)} que organiza as habilidades de alfabetização do 1º ano em uma progressão clara e auditável. \textbf{Importante:} a escala K0--K5 é própria deste livro e não corresponde à nomenclatura oficial do Kumon.

\niveisK

\section{Descrição de Cada Nível}

\subsection{K0 --- Prontidão}

Habilidades preparatórias: coordenação motora fina (pegar o lápis, traçar linhas), atenção sustentada, oralidade (falar e escutar), noção espacial (em cima, embaixo, esquerda, direita) e reconhecimento visual de formas.

\subsection{K1 --- Letras e Grafias}

Reconhecimento das letras nas quatro formas gráficas (imprensa maiúscula e minúscula; cursiva maiúscula e minúscula), traçado correto com direção, nome da letra e associação inicial ao som.

\subsection{K2 --- Consciência Fonológica}

Percepção de sons: rimas, aliterações, contagem de sílabas, identificação do som inicial e final, segmentação e fusão de sílabas (blending e segmenting).

\subsection{K3 --- Decodificação}

Fusão de sílabas para formar palavras (CVC, CCV, CVCV), leitura de palavras simples, correspondência grafema-fonema consistente e primeiras palavras irregulares de alta frequência.

\subsection{K4 --- Frases e Compreensão}

Leitura de frases curtas, compreensão literal, localização de informação explícita, reconhecimento da finalidade de pequenos textos e produção de frases simples.

\subsection{K5 --- Fluência e Transferência}

Leitura com prosódia, textos curtos com compreensão, inferência simples, escrita autônoma de palavras e frases, transferência para situações reais de uso da língua.

\section{Relação entre Níveis K e Unidades}

\begin{center}
\begin{tabular}{|c|l|l|}
\hline
\textbf{Nível K} & \textbf{Habilidade} & \textbf{Unidades do Livro} \\
\hline
K0 & Prontidão & Unidade 0 (novo) \\
\hline
K1 & Letras e Grafias & Unidades 1--3 (vogais e primeiras consoantes) \\
\hline
K2 & Consciência Fonológica & Unidades 4--5 (sílabas diretas e consciência) \\
\hline
K3 & Decodificação & Unidades 6--7 (sílabas e palavras) \\
\hline
K4 & Frases e Compreensão & Unidades 8--9 (frases, nasalizações) \\
\hline
K5 & Fluência e Transferência & Unidade 10 + Parte 3 (Kumon completo) \\
\hline
\end{tabular}
\end{center}

%% ============================================================
%% NOVO CAPÍTULO: ROTAS INDIVIDUALIZADAS (1A-1D)
%% ============================================================

\chapter{Rotas Individualizadas de Aprendizagem (1A--1D)}

\section{Fundamentação}

A Resolução CNE/CEB nº 7/2010, em seus Arts. 29 e 37, determina que o Ensino Fundamental deve atender à \textbf{diversidade de ritmos, condições e necessidades de aprendizagem}. A BNCC (2018) reforça, em suas competências gerais, o respeito à diversidade e a personalização dos percursos formativos.

Com base nesse arcabouço, este livro organiza \textbf{quatro rotas individualizadas (1A a 1D)} que o professor ativa a partir do diagnóstico inicial do aluno.

\rotasindividualizadas

\subsection{Palavra Geradora (inspira\c{c}\~ao em Paulo Freire)}

Freire alfabetizou adultos partindo de \textbf{palavras geradoras} colhidas do
universo vocabular e da viv\^encia dos educandos --- antes de qualquer letra, a
palavra carregava sentido e pertencimento (\textit{``leitura de mundo''}).
Este livro se apropria desse princ\'ipio para a alfabetiza\c{c}\~ao infantil:
cada bloco parte de \textbf{palavras do cotidiano da crian\c{c}a} (casa, m\~ae,
m\~ao, p\~ao, bola, cora\c{c}\~ao), destacadas como \textbf{palavra geradora}
da li\c{c}\~ao. A palavra geradora d\'a o \textbf{sentido}; a t\'ecnica da
Boquinha d\'a o \textbf{som}; a progress\~ao Kumon d\'a a \textbf{ordem e o
dom\'inio}. Ainda em linha com Freire, as fam\'ilias sil\'abicas s\~ao usadas
para a \textbf{recombina\c{c}\~ao} (descobrir novas palavras a partir de
s\'ilabas conhecidas), transformando cada aluno em produtor, n\~ao s\'o
repetidor.

\section{Como Ativar Cada Rota}

\subsection{Rota 1A --- Recomposição Intensiva}

\begin{itemize}
  \item \textbf{Perfil}: aluno com diagnóstico abaixo de 50\% de acertos; não reconhece todas as letras ou não associa letra--som.
  \item \textbf{Estratégia}: retornar aos níveis K0--K1; sessões diárias de 15 minutos; prática multissensorial (boquinha + traçado no ar + massinha); reforço da rotina de 6 passos.
  \item \textbf{Materiais}: Unidade 0 completa + folhas Kumon A-01 a A-05 sem meta de tempo.
\end{itemize}

\subsection{Rota 1B --- Segmentação e Leitura Funcional}

\begin{itemize}
  \item \textbf{Perfil}: diagnóstico entre 50\% e 70\%; reconhece letras, mas falha na fusão de sílabas.
  \item \textbf{Estratégia}: foco em K2--K3; trabalho com sílabas diretas (BA, BE, BI, BO, BU) e palavras CVC; uso sistemático das pistas de boca; leitura de palavras funcionais do cotidiano.
  \item \textbf{Materiais}: Unidades 4--7 + folhas Kumon B-01 a C-03.
\end{itemize}

\subsection{Rota 1C --- Consolidação e Expansão}

\begin{itemize}
  \item \textbf{Perfil}: diagnóstico entre 70\% e 85\%; lê sílabas e palavras simples, mas precisa consolidar frases.
  \item \textbf{Estratégia}: foco em K3--K4; frases curtas; compreensão literal; produção guiada; ampliação de vocabulário.
  \item \textbf{Materiais}: Unidades 8--9 + folhas Kumon C-04 a D-02.
\end{itemize}

\subsection{Rota 1D --- Aprofundamento e Autonomia}

\begin{itemize}
  \item \textbf{Perfil}: diagnóstico acima de 85\%; lê e compreende frases; pronto para fluência.
  \item \textbf{Estratégia}: foco em K4--K5; textos curtos; escrita autônoma; projetos de leitura; transferência para outros componentes curriculares.
  \item \textbf{Materiais}: Unidade 10 + Parte 3 completa (folhas D-03 a D-06) + desafios de escrita.
\end{itemize}

\section{Registro e Acompanhamento das Rotas}

O professor registra a rota de cada aluno na Parte 5 (Guia do Professor) usando o quadro de acompanhamento. A rota não é fixa: a cada reavaliação (mensal ou bimestral), o aluno pode migrar de rota conforme seu progresso. O critério de mudança é o \textbf{domínio demonstrado} (90\% ou mais de acertos com qualidade), nunca a velocidade.

%% ============================================================
%% NOVO CAPÍTULO: DESIGN NEUROINCLUSIVO
%% ============================================================

\chapter{Design Neuroinclusivo do Livro}

\section{Princípios de Acessibilidade para TEA e Dislexia}

O Volume 1 foi reestruturado com base em princípios de design universal para a aprendizagem (DUA) e em evidências de acessibilidade para estudantes neurodivergentes:

\begin{enumerate}[leftmargin=2cm]
  \item \textbf{Rotina visual estável}: todas as lições seguem os mesmos seis passos (observar, ouvir, falar, traçar, ler, conferir), reduzindo a carga de memória de trabalho e a ansiedade de previsibilidade.
  \item \textbf{Instruções curtas}: uma ação principal por bloco; frases de até 12 palavras; vocabulário direto.
  \item \textbf{Alto contraste e espaço em branco}: contraste mínimo WCAG 2.1 AA (4.5:1); margens amplas; baixa densidade visual por página.
  \item \textbf{Pausa e pista explícitas}: cada folha oferece opção de pausa (autorregulação) e pedido de pista (apoio sem estigma).
  \item \textbf{Velocidade não é critério}: o tempo de execução (SCT) é informativo para o professor, mas o único critério de avanço é o domínio demonstrado (90\%+).
  \item \textbf{Multimodalidade}: imagem, som, movimento articulatório e traçado convergem para a mesma unidade fonológica.
\end{enumerate}

\section{A Rotina de Seis Passos}

\rotinavisual{OBSERVAR}{OUVIR}{FALAR}{TRAÇAR}{LER}{CONFERIR}

\begin{description}[leftmargin=2cm]
  \item[1. Observar] A criança olha a letra, a ilustração e a boca do professor.
  \item[2. Ouvir] A criança escuta o som (fonema) isolado e em palavras.
  \item[3. Falar] A criança articula o som, sentindo o movimento da boca.
  \item[4. Traçar] A criança traça a letra nas quatro formas, primeiro no ar, depois no papel.
  \item[5. Ler] A criança lê sílabas e palavras com a letra estudada.
  \item[6. Conferir] A criança confere o próprio trabalho com apoio do professor (autoavaliação).
\end{description}

\section{As Quatro Formas Gráficas das Letras}

Cada lição treina as quatro formas gráficas de cada letra:

\begin{center}
\begin{tabular}{|l|c|c|}
\hline
\textbf{Forma} & \textbf{Exemplo} & \textbf{Uso no livro} \\
\hline
Imprensa maiúscula & \textbf{A} & Cabeçalho, reconhecimento, letreiros \\
\hline
Imprensa minúscula & \textbf{a} & Leitura de textos, palavras \\
\hline
Cursiva maiúscula & \cursiva{A} & Assinatura, cartões, convites \\
\hline
Cursiva minúscula & \cursiva{a} & Escrita pessoal, cópia, ditado \\
\hline
\end{tabular}
\end{center}

\section{A Pista Articulatória (Boquinha)}

A técnica da boquinha é mantida como \textbf{apoio multimodal}, não como recurso obrigatório: a criança que não quiser ou não puder olhar para a boca pode usar as pistas alternativas (toque no pescoço para sentir vibração, espelho, mão à frente da boca para sentir o ar). Cada lição traz a pista articulatória com três dimensões: som, posição da boca e pista tátil.

%% ============================================================

\chapter{Preparação para CAEd e SPAECE}

\section{O que é a Matriz CAEd?}

A Matriz de Referência para Avaliação é desenvolvida pelo Centro de Avaliação Educacional (CAEd) da Universidade Federal de Juiz de Fora (UFJF) e serve de base para a Prova Brasil e o Saeb.

\textbf{Nota importante:} a Matriz CAEd e a Matriz SPAECE-Alfa são \textbf{matrizes avaliativas correlatas} — elas avaliam o que foi ensinado, mas \textbf{não substituem o currículo}. O currículo deste livro é a BNCC (2018). As matrizes são usadas aqui apenas para alinhar as atividades às habilidades avaliadas, permitindo que o estudante se familiarize com o formato das provas.

\textbf{Descritores de Língua Portuguesa:}

\begin{table}[h]
\centering
\caption{Descritores da Matriz CAEd (referência avaliativa correlata)}
\begin{tabular}{|c|p{5cm}|p{4cm}|}
\hline
\textbf{Descritor} & \textbf{Descrição} & \textbf{Avalia} \\
\hline
D1 & Leitura de palavras isoladas & Decodificação \\
\hline
D2 & Leitura de frases & Compreensão literal \\
\hline
D3 & Leitura de textos & Compreensão global \\
\hline
D4 & Produção de texto & Escrita \\
\hline
D5 & Conhecimento gramatical & Metalinguagem \\
\hline
\end{tabular}
\end{table}

\section{SPAECE --- Padrões de Desempenho}

O SPAECE (Sistema Permanente de Avaliação da Educação Básica do Ceará) avalia o desempenho dos alunos em Língua Portuguesa e Matemática. O SPAECE-Alfa, voltado ao 2º ano, verifica o nível de alfabetização dos estudantes do ciclo inicial.

\textbf{Descritores da Matriz SPAECE-Alfa (correlatos ao 1º ano):}

\spaecealfa

\textbf{Níveis de desempenho em LP (1º ano):}

\begin{table}[h]
\centering
\caption{Níveis de Desempenho SPAECE}
\begin{tabular}{|c|p{5cm}|p{5cm}|}
\hline
\textbf{Nível} & \textbf{Descrição} & \textbf{Ação Pedagógica} \\
\hline
I & Não lê nem escreve palavras com sílabas simples & Trabalhar consciência fonológica e sílabas diretas (Rota 1A) \\
\hline
II & Lê e escreve palavras com sílabas diretas & Consolidar sílabas diretas e introduzir indiretas (Rota 1B) \\
\hline
III & Lê e escreve palavras com sílabas diretas e indiretas & Trabalhar sílabas complexas e frases (Rota 1C) \\
\hline
IV & Lê e compreende frases e pequenos textos & Fluência e transferência (Rota 1D) \\
\hline
\end{tabular}
\end{table}

%% ============================================================
%% NOVO CAPÍTULO (R201.3): SEQUÊNCIA DIDÁTICA GRANULAR (LÓGICA KUMON)
%% ============================================================

\chapter{Sequência Didática Granular --- Lógica Kumon}

\section{Princípio da Granularidade}

O método Kumon decompõe cada habilidade em \textbf{micro-passos}: cada folha
introduz \textbf{um único elemento novo} e revisa os anteriores. Este livro
aplica o mesmo princípio em escala sistêmica --- do fonema isolado ao texto
fluente --- com três regras:

\begin{enumerate}[leftmargin=2cm]
  \item \textbf{Um elemento novo por passo}: cada missão/folha acrescenta
        exatamente uma letra, um som ou uma estrutura silábica nova.
  \item \textbf{Revisão cumulativa}: 30\% a 40\% de cada folha retoma passos
        anteriores (revisão integrada nas atividades 1 e 3 de cada missão).
  \item \textbf{Domínio antes do avanço}: só se avança com 90\% ou mais de
        acertos com qualidade, quantas sessões forem necessárias. Velocidade
        nunca é critério.
\end{enumerate}

\section{Mapa de Micro-Passos do Volume 1}

Cada passo é granular: \textbf{novo} = o único elemento inédito; \textbf{revisa}
= o que a folha retoma de passos anteriores.

\begin{center}
\small
\begin{tabular}{|c|l|l|l|}
\hline
\textbf{Passo} & \textbf{Novo elemento (único)} & \textbf{Revisa} & \textbf{Nível K} \\
\hline
1 & Vogal A: 4 formas + som aberto & --- & K1 \\
2 & Vogal E: formas + sons & A & K1 \\
3 & Vogal I: formas + sons & A, E & K1 \\
4 & Vogal O: formas + sons & A, E, I & K1 \\
5 & Vogal U: formas + sons & A, E, I, O & K1 \\
6 & Consoante M: som nasal & vogais & K1 \\
7 & Consoante S: 4 valores & vogais + M & K1 \\
8 & Consoante T & vogais + M, S & K1 \\
9 & Consoante R: forte/fraco & M, S, T & K1 \\
10 & Consoante L & M, S, T, R & K1 \\
11 & Sílabas diretas (CV) com M, S, T, R, L (MA--LA) & M, S, T, R, L & K2 \\
12 & Jogo e frases curtas & sílabas diretas & K2 \\
13 & Consoante B + família & sílabas e vogais & K2 \\
14 & Consoante C: K/S e Ç (cedilha) & B & K2 \\
15 & Consoante D & B, C & K2 \\
16 & Consoante F & D & K2 \\
17 & Consoante G: forte/J & F & K2 \\
18 & Consoante J & G & K2 \\
19 & Consoante K & J & K2 \\
20 & Consoante N: nasal & K & K2 \\
21 & Consoante P & N & K2 \\
22 & Consoante V & P & K2 \\
23 & Consoante Z: Z/S & V & K2 \\
24 & Consoante W & Z & K2 \\
25 & Consoante Y & W & K2 \\
26 & Sílabas nasais: AM, EM, IM, OM, UM / AN, EN, IN, ON, UN & consoantes & K3 \\
27 & Ditongos nasais: ÃO, ÃE, ÕE (e AM no fim de verbos) & nasais & K4 \\
28 & Consoante X: 4 sons & ditongos & K4 \\
29 & Consoante Q: U mudo/sonoro & X & K4 \\
30 & Consoante H: mudo + CH, LH, NH & Q & K4 \\
31 & Fluência e transferência (textos curtos) & tudo & K5 \\
\hline
\end{tabular}
\end{center}

\section{Por que essa ordem (e não a alfabética)?}

\begin{itemize}
  \item \textbf{Vogais primeiro}: são o núcleo de toda sílaba do português.
  \item \textbf{Consoantes regulares e frequentes antes} (M, S, T, R, L):
        permitem formar palavras reais desde cedo (mala, sala, tatu).
  \item \textbf{Letras de múltiplos sons por último} (C, G, X, Q, H):
        exigem regra posicional e vêm depois da base regular consolidada.
  \item \textbf{Sílabas só depois das letras}: fusão (blending) exige
        automatização prévia do grafema--fonema.
  \item \textbf{Nasais e ditongos depois das sílabas diretas}: estruturas
        complexas chegam quando a estrutura CV está dominada.
  \item \textbf{Frases e fluência no fim}: compreensão e prosódia são o teto
        da pirâmide, não a base.
\end{itemize}

\section{Como Usar o Mapa no Dia a Dia}

\begin{enumerate}[leftmargin=2cm]
  \item Marque no mapa o passo atual de cada aluno (o mapa vira painel de
        progresso da turma).
  \item Se o aluno falhar em revisão, \textbf{volte ao passo} onde a falha
        apareceu --- nunca pule.
  \item Use as rotas 1A--1D para decidir a velocidade de percurso, não a
        sequência: a ordem dos passos é igual para todos.
  \item Registre no quadro da Parte 5 a cada mudança de passo.
\end{enumerate}

\checklistdominio{Sequência granular aplicada: 1 novo por passo + revisão cumulativa + domínio 90\%}
